% ============================================================
% EXPERIMENTS — Full version for NeurIPS 2026
% ============================================================

\section{Experiments}
\label{sec:experiments}

We evaluate \methodname{} against four baselines on synthetic and
real-world benchmarks. All code, data generators, and trained models
are available at the URL in the abstract.

\paragraph{Implementation details.}
All experiments use a Mamba model with $d_{\text{model}} = 64$,
$d_{\text{state}} = 16$, $L = 4$ layers, expand factor 2, and causal
convolution kernel size 4. Training uses AdamW (lr = $10^{-3}$,
weight decay = $10^{-4}$) with cosine annealing, for 150 epochs on
subsequences of length 200. No data augmentation or early stopping.
All time series are standardized to zero mean and unit variance before
training. Results are reported over 3 random seeds (mean $\pm$ std
where indicated). All ground truth graphs are defined by the data
generation process (coefficient matrices for VAR, physics equations
for Lorenz, established literature for climate/fMRI). Full
hyperparameters for all methods are in
Appendix~\ref{app:hyperparams}.

\paragraph{Baselines.}
\begin{itemize}[leftmargin=1.5em, topsep=0pt, itemsep=0pt]
    \item \textbf{Granger}~\citep{granger1969investigating}: Bivariate
    F-tests via \texttt{statsmodels}, max lag 5. Score = negative log
    $p$-value.
    \item \textbf{PCMCI}~\citep{runge2019pcmci}: Partial correlation
    conditional independence test, $\tau_{\max} = 5$, via
    \texttt{tigramite}.
    \item \textbf{cLSTM}~\citep{tank2021neural}: Per-target LSTM with
    group-lasso ($\lambda = 0.01$), 200 epochs.
    \item \textbf{Correlation}: Absolute Pearson correlation matrix.
    Non-causal sanity check.
\end{itemize}

\paragraph{Metrics.}
AUROC (primary, threshold-free), AUPRC, F1 at the oracle threshold,
precision, recall, and Spearman rank correlation between predicted and
true causal strengths. We report wall-clock runtime including training.

% ── 4.1 Synthetic ──
\subsection{Synthetic Benchmarks}
\label{sec:synthetic}

We evaluate on five synthetic datasets with known ground-truth causal
graphs (Table~\ref{tab:synthetic}):

\begin{itemize}[leftmargin=1.5em, topsep=0pt, itemsep=0pt]
    \item \textbf{VAR(1) chain K=3, K=5}: Linear chain
    $X_0 \to X_1 \to \cdots \to X_{K-1}$ with coefficient 0.4.
    \item \textbf{VAR(1) fork K=5}: Hub structure with $X_0$ driving all
    others.
    \item \textbf{Regime-switching K=3}: Forward chain for $t < T/2$,
    reversed chain for $t \geq T/2$. Tests non-stationarity robustness.
    \item \textbf{Lorenz K=3}: Nonlinear chaotic system with known
    coupling structure ($X \leftrightarrow Y$, $X \to Z$, $Y \to Z$).
\end{itemize}

\input{table1_synthetic}

\textbf{Key findings.} (1)~\methodname{} achieves perfect AUROC (1.0) on
3 of 5 benchmarks. (2)~It outperforms both PCMCI and cLSTM on
\emph{all 5} datasets. (3)~Granger tests dominate on linear stationary
data---their exact design target---but \methodname{} remains within 0.15
AUROC. (4)~\methodname{} is the only neural method competitive with the
classical gold standard.

% ── 4.2 Non-stationary ──
\subsection{Non-Stationary Data}
\label{sec:nonstationary}

The regime-switching benchmark (Table~\ref{tab:synthetic}, row 4) tests
causal recovery when stationarity is violated. PCMCI's theoretical
guarantees assume stationarity; when this assumption breaks, its
performance should degrade. The selective SSM, by contrast, can adapt its
data-dependent $\bB_t$, $\bC_t$ matrices to regime changes.

% Figure 1 placeholder - will be auto-generated
% \begin{figure}[t]
%     \centering
%     \includegraphics[width=0.9\linewidth]{figures/figure1_windowed_auroc.png}
%     \caption{Windowed AUROC over time on regime-switching data.
%     SSM-Causal adapts to the regime change at $t=2000$, while
%     static methods average over both regimes.}
%     \label{fig:windowed}
% \end{figure}

% ── 4.3 Real-world ──
\subsection{Real-World Data}
\label{sec:realworld}

\paragraph{Climate teleconnections.}
We evaluate on 8 climate indices (ENSO, NAO, PDO, AMO, IOD, PNA, SOI,
SAM) with 720 monthly timesteps (60 years). Ground truth is derived from
established teleconnection relationships in the climate literature
(ENSO $\to$ IOD, ENSO $\to$ PNA, PDO $\to$ ENSO, NAO $\to$ AMO).

\paragraph{Brain functional connectivity.}
We evaluate on 12 brain regions of interest (ROIs) from four functional
networks (default mode, visual, motor, frontoparietal) with 1200
timepoints. Ground truth is from known directed connectivity (V1 $\to$ V2
$\to$ MT, SMA $\to$ M1, mPFC $\leftrightarrow$ PCC).

% Table 3 placeholder
% \input{table3_realworld}

% ── 4.4 Efficiency ──
\subsection{Efficiency Analysis}
\label{sec:efficiency}

A key practical advantage of \methodname{} is computational efficiency.
Once the Mamba model is trained, causal graph extraction requires only
reading two weight matrices and performing a single matrix
multiplication---$O(dK^2)$ operations, typically under 1 millisecond.

We benchmark wall-clock runtime across varying $K$ (number of variables)
and $T$ (sequence length). The training cost is $O(TdNL)$, which is
linear in $T$---compared to PCMCI's $O(T^2)$ scaling from pairwise
conditional independence tests.

% Table 4 placeholder
% \input{table4_efficiency}

% ── 4.5 Ablation ──
\subsection{Ablation Study}
\label{sec:ablation}

\subsection{Why SSM? Architecture Comparison}
\label{sec:architecture}

A natural question is whether the causal score $\bS = |\bW_{\text{out}}
\bW_{\text{in}}|$ works specifically because of the SSM architecture, or
whether any model with input/output projections would suffice. To test
this, we apply the same weight-projection extraction to four architectures
trained on identical data with identical projections: Mamba (SSM), MLP,
LSTM, and a linear model. If \methodname{} outperforms MLP and LSTM on
this extraction, it demonstrates that the SSM's structured state dynamics
help the weight matrices encode causal structure more faithfully.

% TODO: Insert architecture comparison table from exp7 results

\subsection{Ablation Study}
\label{sec:ablation}

We ablate five design choices on VAR(1) chain K=5:

\begin{enumerate}[leftmargin=1.5em, topsep=0pt, itemsep=0pt]
    \item \textbf{Score variant}: Weight-projection vs.\ Jacobian vs.\
    B$\cdot$C$^\top$ projection.
    \item \textbf{Model dimension}: $d \in \{32, 64, 128\}$.
    \item \textbf{State dimension}: $N \in \{8, 16, 32\}$.
    \item \textbf{Number of layers}: $L \in \{2, 4, 6\}$.
    \item \textbf{Data length}: $T \in \{500, 1000, 2000, 5000\}$.
\end{enumerate}

% Table 5 placeholder
% \input{table5_ablation}
